% GENERATED FILE. Edit canonical lessons, then run npm run generate:books.
% canonical-source-hash: fnv1a64:4e60b43c4cb5856f
% canonical-lessons: FR-C67-trouver, FR-C67-poisson, FR-C67-poulet, FR-C67-fruit, FR-C67-pomme

\chapter{Find, Fish, Chicken, Fruit, Apple}
\label{ch:fr-find-fish-chicken-fruit-apple}
\hlchaptermodality{67}

\begin{chapteropening}
\textbf{By the end of this chapter:} \emph{I can name to find, fish, chicken, a fruit, an apple.}

Complete the last lesson of chapter 67: I can name to find, fish, chicken, a fruit, an apple.
\end{chapteropening}

\section[trouver]{trouver — to find}

\label{lesson:FR-C67-trouver}



\begin{quote}
Before the new one: say the French for to win; to earn, then the French for to hide.
\end{quote}

\subsection*{You'll want to know: trouver}
\textbf{trouver} — \textquotedblleft{}to find\textquotedblright{}.

\begin{cousinweb}
\textbf{trouver} $\leftarrow$ Old French \emph{trover}, probably from Vulgar Latin \emph{*tropare}, \textquotedblleft{}to compose a poem\textquotedblright{}, from Latin \emph{tropus}, \textquotedblleft{}a figure of speech\textquotedblright{}, from Greek \emph{tropos}, \textquotedblleft{}a turn\textquotedblright{} — inherited (EtymDB 2.1). This is the usual account; it is not certain.
\end{cousinweb}

\begin{grammarlens}[title={What you've built}]
One more word for this chapter.
\end{grammarlens}

\subsection*{Your turn}
\begin{itemize}
\raggedright
  \item \emph{Say it:} \emph{trouver}
  \item \emph{Say it:} \emph{trouver}, once more
  \item \emph{Say it:} say \emph{trouver}
  \item \emph{Recall:} say the French for to win; to earn, then the French for to hide, then say \emph{trouver} again
\end{itemize}

\begin{tcolorbox}[breakable,colback=teal!4,colframe=teal!35!black,title={Before you move on}]
What does trouver mean? (\textquotedblleft{}To find\textquotedblright{}.) Say it once more.
\end{tcolorbox}
\section[le poisson]{le poisson — fish}

\label{lesson:FR-C67-poisson}



\begin{quote}
Before the new one: say the French for to hide, then the French for to find.
\end{quote}

\subsection*{You'll want to know: le poisson}
\textbf{le poisson} — \textquotedblleft{}fish\textquotedblright{}.

\begin{cousinweb}
\textbf{poisson} $\leftarrow$ Latin \emph{piscis}, \textquotedblleft{}fish\textquotedblright{}, through a Vulgar Latin form \emph{*piscio} — inherited. (EtymDB 2.1)
\end{cousinweb}

\begin{grammarlens}[title={What you've built}]
One more word for this chapter.
\end{grammarlens}

\subsection*{Your turn}
\begin{itemize}
\raggedright
  \item \emph{Say it:} \emph{le poisson}
  \item \emph{Say it:} \emph{le poisson}, once more
  \item \emph{Say it:} say \emph{le poisson}
  \item \emph{Recall:} say the French for to hide, then the French for to find, then say \emph{le poisson} again
\end{itemize}

\begin{tcolorbox}[breakable,colback=teal!4,colframe=teal!35!black,title={Before you move on}]
What does le poisson mean? (\textquotedblleft{}Fish\textquotedblright{}.) Say it once more.
\end{tcolorbox}
\section[le poulet]{le poulet — chicken}

\label{lesson:FR-C67-poulet}



\begin{quote}
Before the new one: say the French for to find, then the French for fish.
\end{quote}

\subsection*{You'll want to know: le poulet}
\textbf{le poulet} — \textquotedblleft{}chicken\textquotedblright{}.

\begin{cousinweb}
\textbf{poulet} — built from the French word for \textquotedblleft{}hen\textquotedblright{}, with \emph{-et}, \textquotedblleft{}little\textquotedblright{}. The hen word $\leftarrow$ Latin \emph{pulla}, the feminine of \emph{pullus}, \textquotedblleft{}a young animal\textquotedblright{} (Lewis \& Short).
\end{cousinweb}

\begin{grammarlens}[title={What you've built}]
One more word for this chapter.
\end{grammarlens}

\subsection*{Your turn}
\begin{itemize}
\raggedright
  \item \emph{Say it:} \emph{le poulet}
  \item \emph{Say it:} \emph{le poulet}, once more
  \item \emph{Say it:} say \emph{le poulet}
  \item \emph{Recall:} say the French for to find, then the French for fish, then say \emph{le poulet} again
\end{itemize}

\begin{tcolorbox}[breakable,colback=teal!4,colframe=teal!35!black,title={Before you move on}]
What does le poulet mean? (\textquotedblleft{}Chicken\textquotedblright{}.) Say it once more.
\end{tcolorbox}
\section[le fruit]{le fruit — a fruit}

\label{lesson:FR-C67-fruit}



\begin{quote}
Before the new one: say the French for fish, then the French for chicken.
\end{quote}

\subsection*{You'll want to know: le fruit}
\textbf{le fruit} — \textquotedblleft{}a fruit\textquotedblright{}.

\begin{cousinweb}
\textbf{fruit} $\leftarrow$ Old French \emph{fruit} $\leftarrow$ Latin \emph{fructus}, \textquotedblleft{}enjoyment, produce\textquotedblright{} — inherited. (EtymDB 2.1)
\end{cousinweb}

\begin{grammarlens}[title={What you've built}]
One more word for this chapter.
\end{grammarlens}

\subsection*{Your turn}
\begin{itemize}
\raggedright
  \item \emph{Say it:} \emph{le fruit}
  \item \emph{Say it:} \emph{le fruit}, once more
  \item \emph{Say it:} say \emph{le fruit}
  \item \emph{Recall:} say the French for fish, then the French for chicken, then say \emph{le fruit} again
\end{itemize}

\begin{tcolorbox}[breakable,colback=teal!4,colframe=teal!35!black,title={Before you move on}]
What does le fruit mean? (\textquotedblleft{}A fruit\textquotedblright{}.) Say it once more.
\end{tcolorbox}
\section[la pomme]{la pomme — an apple}

\label{lesson:FR-C67-pomme}



\begin{quote}
Before the new one: say the French for chicken, then the French for a fruit.
\end{quote}

\subsection*{You'll want to know: la pomme}
\textbf{la pomme} — \textquotedblleft{}an apple\textquotedblright{}.

\begin{cousinweb}
\textbf{pomme} $\leftarrow$ Latin \emph{poma}, \textquotedblleft{}fruits\textquotedblright{}, the plural of \emph{pomum} — inherited. (EtymDB 2.1) In Latin \emph{pomum} was any fruit; French narrowed it to the apple.
\end{cousinweb}

\begin{grammarlens}[title={What you've built}]
One more word for this chapter.
\end{grammarlens}

\subsection*{Your turn}
\begin{itemize}
\raggedright
  \item \emph{Say it:} \emph{la pomme}
  \item \emph{Say it:} \emph{la pomme}, once more
  \item \emph{Say it:} say \emph{la pomme}
  \item \emph{Recall:} say the French for chicken, then the French for a fruit, then say \emph{la pomme} again
\end{itemize}

\begin{tcolorbox}[breakable,colback=teal!4,colframe=teal!35!black,title={Before you move on}]
What does la pomme mean? (\textquotedblleft{}An apple\textquotedblright{}.) Say it once more.
\end{tcolorbox}
